% Part: first-order-logic
% Chapter: models-theories
% Section: theories

\documentclass[../../../include/open-logic-section]{subfiles}

\begin{document}

\olfileid{fol}{mat}{the}

\olsection{প্রথম-ক্রমের তত্ত্বের উদাহরণ}

\begin{ex}
ভাষা~$\Lang L_<$-এ কঠোর রৈখিক ক্রমের তত্ত্বটি নিচের সেট দ্বারা
স্বতঃসিদ্ধায়িত:
\begin{align*}
\{\quad & \lforall[x][\lnot x < x], \\
& \lforall[x][\lforall[y][((x < y \lor y <
    x) \lor x = y)]], \\
& \lforall[x][\lforall[y][\lforall[z][((x < y
      \land y < z) \lif x < z)]]] \quad \}
\end{align*}
এটি অভিপ্রেত !!{structure}s-গুলিকে সম্পূর্ণভাবে ধারণ করে: প্রত্যেক
কঠোর রৈখিক ক্রম এই স্বতঃসিদ্ধ-পদ্ধতির একটি মডেল; এবং বিপরীত দিকে,
$X$ সেটের উপর $R$ একটি কঠোর রৈখিক ক্রম হলে $\Domain M = X$ ও
$\Assign{<}{M} = R$-সহ গঠন~$\Struct M$ এই তত্ত্বের একটি মডেল।
\end{ex}

\begin{ex}
$\Obj 1$ (!!{constant}) ও $\cdot$ (দ্বি-স্থানীয় !!{function})-এর
ভাষায় গোষ্ঠীর তত্ত্বটি নিচের সূত্রগুলি দ্বারা স্বতঃসিদ্ধায়িত:
\begin{align*}
& \lforall[x][\eq[(x \cdot \Obj 1)][x]]\\
& \lforall[x][\lforall[y][\lforall[z][\eq[(x \cdot (y \cdot z))][((x
          \cdot y) \cdot z)]]]]\\
& \lforall[x][\lexists[y][\eq[(x \cdot y)][\Obj 1]]]
\end{align*}
\end{ex}

\begin{ex}
পাটীগণিতের ভাষা~$\Lang L_A$-তে পেয়ানো পাটীগণিতের তত্ত্বটি নিচের
বাক্যগুলি দ্বারা স্বতঃসিদ্ধায়িত:
\begin{align*}
& \lforall[x][\lforall[y][(\eq[x'][y'] \lif \eq[x][y])]]\\
& \lforall[x][\eq/[\Obj 0][x']]\\
& \lforall[x][\eq[(x + \Obj 0)][x]]\\
& \lforall[x][\lforall[y][\eq[(x + y')][(x + y)']]]\\
& \lforall[x][\eq[(x \times \Obj 0)][\Obj 0]]\\
& \lforall[x][\lforall[y][\eq[(x \times y')][((x \times y) + x)]]]\\
& \lforall[x][\lforall[y][(x < y \liff \lexists[z][\eq[(z' + x)][y])]]]\\
\intertext{এবং নিচের রূপের সব বাক্য}
& (!A(\Obj 0) \land \lforall[x][(!A(x) \lif !A(x'))]) \lif \lforall[x][!A(x)]
\end{align*}
শেষের রূপটির অসীমসংখ্যক বাক্য আছে বলে এই স্বতঃসিদ্ধ-পদ্ধতি অসীম।
শেষের রূপটিকে \emph{আরোহ-ছক} বলা হয়। (আসলে আরোহ-ছকটি এখানে যতটা
বলেছি তার চেয়ে একটু বেশি জটিল।)

শেষ স্বতঃসিদ্ধটি~$<$-এর একটি \emph{স্পষ্ট সংজ্ঞা}।
\end{ex}

\begin{ex}
গণিতের ভিত্তিতে (এবং দর্শনেও) বিশুদ্ধ সেটের তত্ত্ব গুরুত্বপূর্ণ ভূমিকা
নেয়। কোনো সেটের সব !!{element}s-ও বিশুদ্ধ সেট হলে সেটটি বিশুদ্ধ।
তাই শূন্য সেটকে বিশুদ্ধ ধরা হয়; কিন্তু কোনো সেটের এমন
!!a{element} থাকলে যা নিজে সেট নয়, সেটটি বিশুদ্ধ নয়। অর্থাৎ, বিশুদ্ধ
সেটগুলি শুধু শূন্য সেট থেকে গঠিত, কোনো “উর-উপাদান” থেকে নয়; উর-উপাদান
হলো এমন বস্তু যা নিজে সেট নয়।

নিচের সূত্রগুলিকে বিশুদ্ধ সেটের একটি তত্ত্বের স্বতঃসিদ্ধ-পদ্ধতি হিসেবে
বিবেচনা করা যেতে পারে:
\begin{align*}
& \lexists[x][\lnot \lexists[y][y \in x]]\\
& \lforall[x][\lforall[y][(\lforall[z][(z \in x \liff z \in y)] \lif
\eq[x][y])]]\\
& \lforall[x][\lforall[y][\lexists[z][\lforall[u][(u \in z \liff
(\eq[u][x] \lor \eq[u][y]))]]]]\\
& \lforall[x][\lexists[y][\lforall[z][(z \in y \liff \lexists[u][(z \in
        u \land u \in x)])]]]\\
\intertext{এবং নিচের রূপের সব বাক্য} &
\lexists[x][\lforall[y][(y \in x \liff !A(y))]]
\end{align*}
% BN-SRC-103 TARGET-CORRECTION-BEGIN
\par\smallskip\noindent\textnormal{\small\textbf{উৎস-সংশোধন (BN-SRC-103):} হিমায়িত ইংরেজি উৎসে দ্বিতীয় স্বতঃসিদ্ধের অন্তঃস্থ সার্বিক পরিমাণসূচকে সূত্র-আর্গুমেন্টের উদ্বোধনী বর্গবন্ধনী অনুপস্থিত ছিল: \texttt{\string\lforall[z](...)}। বাংলা পাঠে অন্য পরিমাণসূচকগুলির মতো \texttt{\string\lforall[z][(...)]} পুনরুদ্ধার করা হয়েছে।} % BN-SRC-103 TARGET-CORRECTION-END
প্রথম স্বতঃসিদ্ধটি বলে, কোনো !!{element}s-বিহীন একটি সেট আছে (অর্থাৎ
$\emptyset$ আছে); দ্বিতীয়টি সেটের সদস্যভিত্তিক সমতার নীতি দেয়; তৃতীয়টি
বলে যেকোনো সেট $X$ ও~$Y$-এর জন্য $\{X,Y\}$ সেটটি আছে; এবং চতুর্থটি
বলে যেকোনো সেট~$X$-এর জন্য $\cup X$ আছে, যেখানে $\cup X$ হলো
$X$-এর সব উপাদানের সংঘ।

সব শেষে উল্লিখিত !!{sentence}s-গুলিকে একত্রে
\emph{অনানুষ্ঠানিক ধর্মনির্দেশে সেট-গঠন ছক} বলা হয়। মূলত এটি বলে যে,
প্রত্যেক~$!A(x)$-এর জন্য সেট~$\Setabs{x}{!A(x)}$ আছে—তাই প্রথম দর্শনে
এটি সত্য, উপকারী, এমনকি হয়তো প্রয়োজনীয় স্বতঃসিদ্ধও মনে হয়।
একে “অনানুষ্ঠানিক” বলা হয়, কারণ শেষ পর্যন্ত দেখা যায় এটি তত্ত্বটিকে
অপরিতৃপ্তিযোগ্য করে: $!A(y)$ হিসেবে $\lnot y \in y$ নিলে
!!{sentence}
\[
\lexists[x][\lforall[y][(y \in x \liff \lnot y \in y)]]
\]
পাওয়া যায়, এবং কোনো !!{structure}-এই !!{sentence} পরিতৃপ্ত হয় না।
\end{ex}

\begin{ex}
\emph{অংশতত্ত্বে} \emph{অংশ-সম্পর্ক} একটি মৌলিক সম্পর্ক। সেটের তত্ত্বের
মতোই অংশ-সম্পর্কেরও নানা তত্ত্ব আছে, যেগুলি এই সম্পর্কের বিভিন্ন
(কখনও পরস্পরবিরোধী) ধারণাকে স্বতঃসিদ্ধায়িত করে।

অংশতত্ত্বের ভাষায় একটিমাত্র দ্বি-স্থানীয় বিধেয়
সংকেত~$\Obj P$ থাকে, এবং $\Atom{\Obj P}{x, y}$-এর “অর্থ” হলো $x$,
$y$-এর একটি অংশ। এই ব্যাখ্যা মনে রাখা হলে ওই ভাষার
!!a{structure}-কে \emph{অংশ-গঠন} বলা হয়। অবশ্য একটিমাত্র দ্বি-স্থানীয়
বিধেয়ের প্রত্যেক গঠনই সত্যিকার অর্থে এই নামের যোগ্য নয়।
“অংশ-সম্পর্ক” ধারণ করার সম্ভাবনা থাকতে হলে~$\Assign{\Obj P}{M}$-কে
কিছু শর্ত পরিতৃপ্ত করতে হবে; অংশ-সম্পর্কের তত্ত্বে আমরা সেগুলিকে
স্বতঃসিদ্ধ হিসেবে রাখতে পারি। যেমন, বস্তুগুলির উপর অংশ-সম্পর্ক একটি
আংশিক ক্রম: প্রত্যেক বস্তু নিজেরই একটি অংশ (যদিও
\emph{অযথার্থ} অংশ); দুটি ভিন্ন বস্তু পরস্পরের অংশ হতে পারে না; আর
কোনো বস্তুর অংশের একটি অংশ নিজেও সেই বস্তুর অংশ। লক্ষ করো, এই অর্থে
“এর অংশ” সম্পর্কটি “এর উপসেট” সম্পর্কের অনুরূপ, কিন্তু “এর উপাদান”
সম্পর্কের অনুরূপ নয়—শেষের সম্পর্কটি প্রতিবিম্বও নয়, পরিযায়ীও নয়।
\begin{align*}
& \lforall[x][\Atom{\Obj P}{x,x}] \\
& \lforall[x][\lforall[y][((\Part{x}{y} \land \Part{y}{x})
      \lif \eq[x][y])]] \\
& \lforall[x][\lforall[y][\lforall[z][((\Part{x}{y} \land
        \Part{y}{z}) \lif \Part{x}{z})]]]\\
\intertext{তদুপরি, যেকোনো দুটি বস্তুর একটি অংশতাত্ত্বিক যোগ আছে
  (অর্থাৎ, এমন একটি বস্তু যার অংশ ওই দুটি বস্তু এবং যা এই ধর্মের
  দিক থেকে ন্যূনতম)।}  &
\lforall[x][\lforall[y][\lexists[z][\lforall[u][(\Part{z}{u} \liff
        (\Part{x}{u} \land \Part{y}{u}))]]]]
\end{align*}
এগুলি অধিবিদ্যাবিদদের বিবেচিত অংশ-সম্পর্কের মৌলিক নীতিগুলির মাত্র কয়েকটি।
কিন্তু কিছু সংজ্ঞায়িত সম্পর্ক আগে প্রবর্তন না করলে আরও নীতি সূত্রবদ্ধ
করা বা লিখে ফেলা দ্রুতই কঠিন হয়ে ওঠে। যেমন, অংশতত্ত্বে আগ্রহী অধিকাংশ
অধিবিদ্যাবিদ নিচের নীতিটিকেও বৈধ মনে করেন: কোনো বস্তু~$x$-এর যথার্থ
অংশ~$y$ থাকলেই তার এমন একটি অংশ~$z$-ও থাকে যার সঙ্গে~$y$-এর কোনো
অংশ সাধারণ নয়, এবং $y$ ও~$z$-এর সংযোজন হলো~$x$।
\end{ex}

\end{document}
